\phantomsection\label{main}
\section{Declaration of Scranton}\label{declaration-of-scranton}

A Profession of Faith and Declaration formulated by the Polish National
Catholic Bishops Assembled at Lancaster, New York April 28, 2008

We faithfully adhere to the Rule of Faith laid down by St. Vincent of
Lerins in these terms: “\emph{Id teneamus, ubique, quod semper, quod
ab omnibus creditum est; hoc est etenim vere proprieque catholicum.}�
(We hold that which has been believed everywhere, always, and of all
people; for that is truly and properly Catholic.) For this reason we
persevere in professing the faith of the primitive Church, as formulated
in the ecumenical symbols and specified precisely by the unanimously
accepted decisions of the Ecumenical Councils held in the undivided
Church of the first thousand years.

Therefore, we reject the innovations of the First Vatican Council that
on July 18, 1870 promulgated the dogma of papal infallibility and the
universal Episcopate of the Bishop of Rome, which contradict the Faith
of the ancient Church and which destroy its ancient canonical
constitution by attributing to the Pope the plenitude of ecclesiastical
powers over all dioceses and over all the faithful. By denial of his
primatial jurisdiction we do not wish to deny the historic primacy which
several Ecumenical Councils and the Fathers of the ancient Church have
attributed to the Bishop of Rome by recognizing him as the Primus inter
pares (first among equals).

We also reject the dogma of the Immaculate Conception promulgated by
Pius IX in 1854 in defiance of the Holy Scriptures and in contradiction
to the Tradition of the first centuries.

We further reject the dogmatization of the Catholic teaching of the
bodily Assumption of the Blessed Virgin Mary by Pius XII in 1950 as
being in defiance of the Holy Scriptures.

We reject the contemporary innovations promulgated by the Anglican
Communion and the Old Catholic Churches of the Union of Utrecht. We also
regard these innovations as being in defiance of the Holy Scriptures and
in contradiction to the Tradition of the first centuries, namely: the
ordination of women to the Holy Priesthood, the consecration of women to
the Episcopate and the blessing of same-sex unions.

Considering that the Holy Eucharist (Holy Mass) has always been the
central point of Catholic worship, we consider it our duty to declare
that we maintain with perfect fidelity the ancient Catholic doctrine
concerning the Sacrament of the Altar, by believing that we receive the
Body and the Blood of our Savior Jesus Christ under the species of bread
and wine. The Eucharistic celebration in the Church is neither a
continual repetition nor a renewal of the expiatory sacrifice which
Jesus offered once for all upon the Cross, but it is a sacrifice because
it is the perpetual commemoration of the sacrifice offered upon the
Cross; and it is the act by which we represent upon earth and
appropriate to ourselves the one offering which Jesus Christ makes in
Heaven, according to the Epistle to the Hebrews 9:11,12, for the
salvation of redeemed humanity, by appearing for us in the presence of
God (Hebrews 9:24). The character of the Holy Eucharist being thus
understood, it is at the same time, a sacrificial feast by means of
which the faithful in receiving the Body and Blood of our Savior enter
into communion with one another (1 Corinthians 10:17).

We hope that Catholic theologians, by maintaining the faith of the
undivided Church, will succeed in establishing an agreement in regard to
all such questions that have caused controversy ever since the Church
became divided.

We exhort the priests under our jurisdiction: to teach the essential
Christian truths by the proclamation of the Word of God and by the
instruction of the faithful; to seek and practice charity when
discussing controversial doctrines; and in word and deed to set, in
accordance with the teachings of our Savior Jesus Christ, an example for
the faithful of the Church.

By faithfully maintaining and professing the doctrine of Jesus Christ,
by refusing to accept those errors that have crept into the Church by
human fault, and by repudiating the abuses in ecclesiastical matters and
the tendency of some Church leaders to seek temporal wealth and power,
we believe that we will effectively combat the great evils of our day,
which are unbelief and indifference in matters of faith.

Most Rev. Robert M. Nemkovich\\
Rt. Rev. Thomas J. Gnat\\
Rt. Rev. Thaddeus S. Peplowski\\
Rt. Rev. Jan Dawidziuk\\
Rt. Rev. Sylvester Bigaj\\
Rt. Rev. Anthony Mikovsky\\
Rt. Rev. Anthony D. Kopka\\
Rt. Rev. John E. Mack
